\bookchapter{פרק 65: פרויקט 1 — ספר החשבונות שלי (גרסה מקומית)}

\chapterintro{ברוכים הבאים ל“פרויקט הרציני” הראשון שלכם. ספר החשבונות הזה לא מתחבר לרשת, אבל יש בו הכול: אפשר לרשום תנועה, לצפות בפירוט, לסכם לפי חודש ולמחוק ולערוך. אחרי שתסיימו אותו, תעברו מסע מלא של “נתונים לוגיקה ממשק”. אל תמהרו להעתיק קוד — קודם חושבים לעומק ואז עושים; זה העיקרון הקבוע שלנו.}

\booksection{חשיבה על דרישות: קודם שואלים חמש שאלות}

\begin{center}
\begin{tabularx}{\linewidth}{|X|X|}
\hline
\textbf{מרכיב} & \textbf{תשובה} \tabularnewline
\hline
Who (למי זה מיועד) & לשימוש עצמי; גם משפחה וחברים \tabularnewline
What (מה עושים) & רישום הכנסות והוצאות, צפייה בפירוט התנועות, סיכום חודשי, סטטיסטיקה לפי קטגוריה \tabularnewline
Why (למה) & תרגול ידיים + הבנה לאן הכסף הולך \tabularnewline
How (איך) & Swift + SwiftUI + SwiftData + MVVM, מקומי לחלוטין בלי שרת \tabularnewline
Constraint (אילוצים) & להשלים תוך שבוע; תמיכה מ-iOS 17 ומעלה (מחסום SwiftData) \tabularnewline
\hline
\end{tabularx}
\end{center}

\begin{center}\small\color{gray}חמשת מרכיבי הדרישות של ספר החשבונות\end{center}

ואז מפרקים אותו לשלבים קטנים ש“כל אחד מהם שמיש”: (1) בניית שכבת נתונים ב-SwiftData (שתדע לשמור רשומה אחת); (2) מסך רשימה (שיציג את התנועות); (3) חלון הזנה (שיאפשר להוסיף רשומה); (4) מחיקה ועריכה; (5) כרטיס סיכום חודשי. בסיום כל שלב, האפליקציה היא “משהו שאפשר להשתמש בו”.

\booksection{שכבת הנתונים: נותנים ל“רישומים” בית}

קודם מגדירים איך נראית רשומה אחת — השדות צריכים להספיק לצרכים, אבל אל תהיו חמדנים:

\begin{lstlisting}[language=Swift]
import SwiftData

@Model
class Record {
    var type: String      // "income" הכנסה | "expense" הוצאה
    var amount: Double    // סכום (יואן)
    var category: String  // קטגוריה: אוכל, תחבורה, משכורת…
    var note: String      // הערה (רשות)
    var date: Date        // תאריך הרישום

    init(type: String, amount: Double, category: String, note: String = "", date: Date = Date()) {
        self.type = type
        self.amount = amount
        self.category = category
        self.note = note
        self.date = date
    }
}
\end{lstlisting}

\begin{tcbinfo}{טיפ: למה הסכום מסוג Double?}
אפליקציות פיננסיות רציניות משתמשות ב-Decimal (עשרון מדויק) כדי להימנע משגיאות נקודה צפה (הבעיה הישנה והמפורסמת שבה 0.1+0.2 $\neq$ 0.3). לפרויקט אימון Double מספיק בהחלט; אם יום אחד תפתחו מוצר פיננסי אמיתי, זכרו לעבור ל-Decimal או ל-“מספר Int שמחושב באגורות” — מלכודת שמפתחים רבים לפניכם נפלו בה.

\end{tcbinfo}

\booksection{MVVM: מפרידים בין הממשק ללוגיקה}

לתת לממשק לגשת ישירות למסד הנתונים? זה יעבוד, אבל קוד הממשק יהפוך ארוך ומבולגן. ארכיטקטורת MVVM מחלקת תפקידים: View (הממשק) רק מצייר, ViewModel (מרכז הלוגיקה) מנהל נתונים, ו-Model (SwiftData) מנהל אחסון.

תצוגה (ממשק SwiftUI)

ViewModel (מוקד הלוגיקה @Observable)

Model (שכבת נתונים SwiftData)

\begin{center}\small\color{gray}איור 4: מבנה שלוש השכבות של MVVM — הנתונים זורמים בכיוון אחד, וכל שכבה נשארת בגבולותיה\end{center}

\begin{lstlisting}[language=Swift]
import Observation

@Observable
class RecordViewModel {
    var records: [Record] = []
    private let context: ModelContext

    init(context: ModelContext) { self.context = context }

    func load() {
        // סינון עם #Predicate של SwiftData + מיון
        let descriptor = FetchDescriptor<Record>(
            sortBy: [SortDescriptor(\.date, order: .reverse)]
        )
        records = (try? context.fetch(descriptor)) ?? []
    }

    func add(type: String, amount: Double, category: String, note: String) {
        let r = Record(type: type, amount: amount, category: category, note: note)
        context.insert(r)
        try? context.save()
        load()                       // טעינה מחדש, הממשק מתרענן
    }

    func delete(_ record: Record) {
        context.delete(record)
        try? context.save()
        load()
    }
}
\end{lstlisting}

בממשק כותבים @State private var vm = RecordViewModel(...) ואז מציגים את הרשימה ישירות מ-vm.records — ברגע שהנתונים משתנים, הממשק מתרענן אוטומטית.

\booksection{הממשק: מסך רשימה + חלון הזנה}

\begin{lstlisting}[language=Swift]
struct RecordListView: View {
    @State private var vm: RecordViewModel
    @State private var showAdd = false

    var body: some View {
        NavigationStack {
            List {
                Section("יתרת החודש") {          // כרטיס סיכום
                    SummaryRow(records: vm.records)
                }
                Section("רשומות") {
                    ForEach(vm.records) { r in
                        RecordRow(record: r)
                    }
                    .onDelete { idx in
                        idx.forEach { vm.delete(vm.records[$0]) }
                    }
                }
            }
            .navigationTitle("פנקס ההוצאות שלי")
            .toolbar {
                Button { showAdd = true } label: {
                    Image(systemName: "plus")
                }
            }
            .sheet(isPresented: $showAdd) {
                AddRecordView(vm: vm)          // דיאלוג הזנה
            }
        }
    }
}
\end{lstlisting}

עיקרי חלון ההזנה: לשדה הסכום מגדירים מקלדת .keyboardType(.decimalPad) (מקלדת מספרים), להכנסה/הוצאה — פקד מקטעים (Picker), לקטגוריה — בחירה מתפריט, ובלחיצה על “שמור” קוראים ל-vm.add(...) וסוגרים עם dismiss().

\begin{tcbwarn}{שימו לב: אל תתעצלו על בדיקות קלט}
חובה לבדוק את שדה הסכום: ערך ריק, מספר שלילי וקלט בלתי חוקי כמו “abc” — את כולם צריך לחסום (בודקים עם Double(text); אם ההמרה נכשלת מתקבל nil). אחת הסיבות הנפוצות ביותר לקריסת אפליקציות היא “להניח שהקלט של המשתמש תמיד מתנהג יפה”. זכרו: כל מה שמגיע מהמשתמש בודקים כאילו זה אויב.

\end{tcbwarn}

\booksection{עריכה, מחיקה וסיכום}

\begin{itemize}
  \item מחיקה: .onDelete של List נותן מחיקה בהחלקה שמאלה + חלון אישור (מחיקה בטעות היא הדבר שמשתמשי אפליקציות ניהול חשבונות שונאים יותר מכל);
  \item עריכה: לחיצה על שורה פותחת חלון עריכה עם הערכים המקוריים ממולאים מראש (לתקן מילה אחת קל יותר מלהקליד הכל מחדש);
  \item סיכום: מחשבים "הכנסות / הוצאות / יתרה של החודש" באמצעות filter + reduce ומציגים בכרטיס שבראש המסך:
\end{itemize}

\begin{lstlisting}[language=Swift]
func monthSummary(_ records: [Record]) -> (income: Double, expense: Double) {
    let cal = Calendar.current
    let monthRecords = records.filter {
        cal.isDate($0.date, equalTo: Date(), toGranularity: .month)
    }
    let income = monthRecords
        .filter { $0.type == "income" }
        .reduce(0) { $0 + $1.amount }
    let expense = monthRecords
        .filter { $0.type == "expense" }
        .reduce(0) { $0 + $1.amount }
    return (income, expense)
}
\end{lstlisting}

\booksection{קבלה: שלוש שכבות בדיקה + שני אתגרים למתקדמים}

אחרי שסיימתם את הפרויקט, ערכו קבלה בשלוש שכבות לפי הסדר “תפקוד קצוות ארכיטקטורה”: האם אפשר לרשום, האם אפשר למחוק, והאם הנתונים שורדים הפעלה מחדש של האפליקציה (תפקוד); האם הזנת 0 או מספר שלילי גורמת לקריסה (קצוות); האם הממשק והנתונים מעורבבים זה בזה, והאם יש ערכים מקודדים קשיח בשדות (ארכיטקטורה).

אחרי שהגרסה הבסיסית עברה קבלה, תנו לעצמכם שתי “מנות מתקדמות”:

\begin{enumerate}
  \item סטטיסטיקה לפי קטגוריות: מוסיפים כרטיסיית "סטטיסטיקה" ומציירים את ההוצאות לפי קטגוריה כתרשים עוגה/עמודות עם Swift Charts — ספריית התרשימים הרשמית של אפל; כמה שורות קוד ויש גרף;
  \item סנכרון iCloud: הפעלת סנכרון iCloud ב-SwiftData דורשת רק כמה הגדרות — ספר החשבונות מסתנכרן אוטומטית בין iPhone ל-iPad, ופתאום הכל נראה "מתקדם" באמת.
\end{enumerate}

\begin{tcbpractice}{הידיים על ההגה: מסיימים את זה השבוע}
עבדו לפי הסדר “שכבת נתונים ViewModel מסך רשימה חלון הזנה כרטיס סיכום”, שלב אחד בכל יום. נתקעתם? שלחו את הודעת השגיאה כמו שהיא למנוע חיפוש או ל-AI (זוכרים את החלק השביעי “פרקטיקת תכנות בעזרת AI”? בתיאור הבעיה הקפידו לציין “התנהגות צפויה, התנהגות בפועל ושלבי שחזור”). אחרי שתסיימו, נהלו איתו את החשבונות שבוע — כשמשתמשים באפליקציה שבניתם בעצמכם, תחושת ההישג דוחפת אתכם להמשיך קדימה.

\end{tcbpractice}

\begin{tcbsummary}{סיכום הפרק}
\begin{itemize}
  \item חמישה מרכיבי דרישות פירוק לשלבים קטנים ש"כל אחד שמיש", ורק אז מתחילים לעבוד
  \item שכבת נתונים: מודל @Model + מיכל modelContainer + שאילתות FetchDescriptor
  \item חלוקת תפקידים ב-MVVM: View רק מצייר, ViewModel עם @Observable מנהל את הלוגיקה, ו-SwiftData מנהל את האחסון
  \item ממשק: List + Section + toolbar + חלון הזנה עם sheet
  \item קבלה בשלוש שכבות (תפקוד/קצוות/ארכיטקטורה) + אתגרים למתקדמים (Swift Charts, סנכרון iCloud)
\end{itemize}

\end{tcbsummary}
